% =====================================================================
%  Information and Coding 2026/27 — Lab work n.º 1
%  Relatório 1: descrição do trabalho (máximo 12 páginas)
%  Nota: o enunciado proíbe código no relatório; usar diagramas,
%  tabelas e gráficos.
% =====================================================================
\documentclass[a4paper,11pt]{article}

\usepackage[utf8]{inputenc}
\usepackage[T1]{fontenc}
\usepackage[portuguese]{babel}
\usepackage[margin=2.2cm]{geometry}
\usepackage{graphicx}
\usepackage{booktabs}
\usepackage{amsmath,amssymb}
\usepackage{siunitx}
\usepackage{xcolor}
\usepackage{subcaption}
\usepackage{hyperref}
\hypersetup{colorlinks=true,linkcolor=blue!50!black,urlcolor=blue!50!black,citecolor=blue!50!black}

% Marcador para texto por escrever (remover antes da entrega)
\newcommand{\todo}[1]{\textcolor{red}{\textbf{[TODO:} #1\textbf{]}}}
\newcommand{\prog}[1]{\texttt{#1}}

\title{Information and Coding --- Trabalho Prático n.º 1\\[4pt]
\large Relatório 1: Ferramentas de áudio e codecs com e sem perdas}
\author{%
  Nome Apelido 1 (n.º mec.~XXXXX) \and
  Nome Apelido 2 (n.º mec.~XXXXX) \and
  Nome Apelido 3 (n.º mec.~XXXXX)}
\date{Universidade de Aveiro --- DETI\\ Outubro de 2026}

\begin{document}
\maketitle

\begin{abstract}
\todo{Para este programa, foi utilizada a linguagem de programação C++ para programas e codecs, uma vez que, tanto o tempo de execução como a memória são mais eficientes. Foi também utilizada a linguagem de programação python para a visualização de gráficos e benchamrks, uma vez que não interfere no custo de desempenho. FALAR AINDA DE preditor, codificação entrópica, DCT) e os principais resultados
(taxa de compressão face ao FLAC, SNR face ao MP3/AAC).}
\end{abstract}

\tableofcontents
% (Remover o índice se for preciso poupar espaço para caber nas 12 páginas)

% =====================================================================
\section{Introdução}
% =====================================================================
\subsection{Objetivos}
{O objetivo da parte I deste projeto é criar um grupo de programas que analisem, quantizem, comparem e aplicem efeitos a ficheiros de áudio WAV, através de histogramas, redução de bits, métricas de erro/SNR e efeitos sonoros como ecos e modulação de amplitude. Na parte II, foi desenvolvido  um codec de áudio lossless que comprima e descomprima ficheiros WAV exatamente, otimizando compressão, velocidade e memória, e compará-lo com codecs como FLAC, Zstd e LZMA. Na parte III foi desenvolvido um codec de áudio lossy baseado em DCT, otimizado para a melhor relação taxa-distorção, justificando experimentalmente as decisões de projeto e comparando-o com codecs MPEG em compressão, SNR, velocidade e memória}

\subsection{Organização do relatório}
A Secção~2 descreve a arquitetura do sistema e as escolhas tecnológicas:
as linguagens de programação, as bibliotecas utilizadas, a organização do
repositório e o conjunto de ficheiros de áudio usado nos testes. As Secções~3,
4 e~5 correspondem às três partes do trabalho, as ferramentas de análise e
manipulação de áudio, o codec sem perdas e o codec com perdas, apresentando,
para cada uma, as decisões tomadas, a sua justificação e os resultados obtidos.
A Secção~6 reúne as conclusões e o trabalho futuro, e a Secção~7 contém um
manual de instalação e utilização do software.
% =====================================================================
\section{Arquitetura e escolhas tecnológicas}
% =====================================================================
\subsection{Linguagens de programação}
Os programas de processamento e os codecs operam amostra a amostra, percorrendo mais de dois milhões de amostras por ficheiro, e o seu desempenho, em tempo e em memória, é determinante quando comparados com ferramentas escritas em C, como o FLAC e o zstd. Por esse motivo foram implementados em C++, que oferece ainda o controlo ao nível do bit e do inteiro exigido pela escrita bit a bit e pelos códigos de Golomb, e integra diretamente a libsndfile. O Python foi reservado para as tarefas em que o custo computacional é irrelevante: a produção dos gráficos, com o matplotlib, os scripts de medição de desempenho e a verificação independente dos resultados através do numpy. Em síntese, C++ onde o custo é por amostra, Python onde o custo é por experiência

\subsection{Bibliotecas e dependências}
\todo{libsndfile (I/O de WAV), numpy/matplotlib, outras.}

\subsection{Organização do repositório e compilação}
\subsection{Organização do repositório e compilação}
O repositório separa a base comum, em \texttt{src/common/}, dos programas que a
utilizam, cada um num ficheiro \texttt{src/nome.cpp}, ficando os scripts de
visualização e de medição em \texttt{scripts/}, a documentação em \texttt{docs/}
e os ficheiros de áudio em \texttt{data/}, estes últimos excluídos do controlo
de versões juntamente com os artefactos de compilação. A base comum reúne os
módulos usados por todos os programas: \texttt{wav\_io.hpp}, responsável pela
leitura e escrita de ficheiros WAV sobre a libsndfile, e \texttt{channels.hpp},
que implementa as transformações MID/SIDE.
\todo{Acrescentar, quando existirem: bitstream.hpp (escrita e leitura bit a bit),
golomb.hpp (codificação entrópica) e metrics.hpp (MSE, L-infinito e SNR,
partilhadas pelo wav\_cmp e pelo codec com perdas).}

A compilação é gerida por CMake, que deteta automaticamente os ficheiros em
\texttt{src/} e gera um executável por programa, ligado à libsndfile através do
\texttt{pkg-config}; acrescentar um programa não exige alterar o ficheiro de
configuração, o que evita conflitos no trabalho em paralelo. Utiliza-se C++20,
indispensável para garantir o comportamento do deslocamento à direita em valores
negativos, em modo \texttt{Release} com otimização \texttt{-O2} e com os avisos
do compilador ativos.

\subsection{Conjunto de dados de teste}
Os testes foram realizados com os sete ficheiros disponibilizados no eLearning,
resumidos na Tabela~\ref{tab:dataset}. Todos partilham o mesmo formato, o que
isola o efeito do conteúdo nos resultados obtidos.
\todo{Uma frase sobre a diversidade de conteúdos e, se for o caso, sobre o
sample07 estar saturado (picos em $\pm$32767, visíveis no histograma).}

\begin{table}[h]
\centering
\caption{Conjunto de ficheiros de áudio usado nos testes. Todos em formato WAV,
PCM inteiro de \SI{16}{bits}, estéreo e \SI{44,1}{\kilo\hertz}.}
\label{tab:dataset}
\sisetup{group-separator={\,}}
\begin{tabular}{l S[table-format=2.1] S[table-format=7.0] S[table-format=1.1] l}
\toprule
Ficheiro & {Duração (s)} & {Amostras/canal} & {Tamanho (MiB)} \\
\midrule
\texttt{sample01.wav} & 29.3 & 1294188 & 4.9\\
\texttt{sample02.wav} & 14.7 &  647388 & 2.5\\
\texttt{sample03.wav} & 20.0 &  882588 & 3.4\\
\texttt{sample04.wav} & 13.3 &  588588 & 2.2\\
\texttt{sample05.wav} & 20.7 &  911988 & 3.5\\
\texttt{sample06.wav} & 24.0 & 1058988 & 4.0\\
\texttt{sample07.wav} & 21.3 &  941388 & 3.6\\
\bottomrule
\end{tabular}
\end{table}

% =====================================================================
\section{Parte I --- Ferramentas de processamento de áudio}
% =====================================================================
\subsection{\prog{wav\_hist}: histogramas}
\subsubsection{Histograma por canal}
\todo{Método e exemplo de gráfico (L e R).}
\subsubsection{Canais MID e SIDE}
\todo{Definição com divisão inteira; convenção de arredondamento adotada.}
\subsubsection{Reconstrução exata de L e R a partir de MID e SIDE}
\todo{Resposta à pergunta da nota de rodapé 2 do enunciado.}
\subsubsection{Bins mais grosseiros ($2^k$ valores)}
\todo{Implementação e comparação visual para vários $k$.}
\subsubsection{Análise dos resultados}
\todo{O que os histogramas mostram (distribuição de L/R vs.\ SIDE, entropia)
e a relação com a codificação das Partes II e III.}

\subsection{\prog{wav\_quant}: quantização escalar uniforme}
\todo{Método (redução de bits, arredondamento vs.\ truncatura), formato de saída.}

\subsection{\prog{wav\_cmp}: métricas de erro}
\subsubsection{Métricas implementadas}
\todo{Definições de $L^2$ (MSE), $L^\infty$ e SNR, por canal e para a média.}
\begin{equation}
  \mathrm{SNR} = 10\log_{10}\frac{\sum_n x[n]^2}{\sum_n \left(x[n]-\tilde{x}[n]\right)^2}\ \text{dB}
\end{equation}
\subsubsection{Resultados da quantização}
\todo{Tabela/gráfico: SNR e $L^\infty$ em função do n.º de bits
(p.\,ex.\ 16, 12, 8, 6, 4, 2). Comparar com o valor teórico de $\approx$\,6{,}02\,dB/bit.}

\subsection{\prog{wav\_effects}: efeitos de áudio}
\subsubsection{Eco simples e ecos múltiplos}
\todo{Equações às diferenças e parâmetros.}
\subsubsection{Modulação de amplitude}
\todo{}
\subsubsection{Atrasos variáveis no tempo}
\todo{}
\subsubsection{Outros efeitos}
\todo{Opcional.}

% =====================================================================
\section{Parte II --- Codec de áudio sem perdas}
% =====================================================================
\subsection{Visão geral e diagrama de blocos}
\todo{Figura com o diagrama de blocos do codificador e do descodificador.}

\subsection{Formato do ficheiro comprimido}
\todo{Cabeçalho (metadados do WAV necessários para reconstrução bit a bit),
organização em blocos/\emph{frames}.}

\subsection{Módulo \emph{BitStream}}
\todo{Escrita/leitura bit a bit, eficiência (\emph{buffering}).}

\subsection{Descorrelação entre canais}
\todo{MID/SIDE e alternativas; escolha adaptativa por bloco.}

\subsection{Predição temporal}
\todo{Preditores polinomiais fixos de ordem 0--3 e/ou LPC; escolha por bloco;
justificação experimental.}

\subsection{Codificação entrópica dos resíduos}
\subsubsection{Códigos de Golomb/Rice}
\todo{Mapeamento de inteiros com sinal, escolha do parâmetro $m$.}
\subsubsection{Adaptação do parâmetro}
\todo{Estimação por bloco vs.\ adaptação contínua.}

\subsection{Parâmetros e opções do codificador}
\todo{Tamanho de bloco, ordem do preditor, etc.; efeito de cada opção.}

\subsection{Avaliação experimental}
\subsubsection{Metodologia}
\todo{Máquina, compilador e flags, medição do tempo e da memória
(\prog{/usr/bin/time -v}), n.º de repetições.}
\subsubsection{Estudo dos parâmetros do codec}
\todo{Gráficos: taxa de compressão vs.\ tamanho de bloco, ordem, etc.}
\subsubsection{Comparação com o estado da arte}
\todo{FLAC (vários níveis), zstd, lzma/xz, gzip, bzip2...}

\begin{table}[h]
\centering
\caption{Comparação de codecs sem perdas (médias sobre o conjunto de teste).}
\label{tab:lossless}
\begin{tabular}{lcccc}
\toprule
Codec & Taxa de compressão & T.\ codif.\ (s) & T.\ descodif.\ (s) & Memória (MB) \\
\midrule
O nosso & -- & -- & -- & -- \\
FLAC -8 & -- & -- & -- & -- \\
zstd -19 & -- & -- & -- & -- \\
xz -9 & -- & -- & -- & -- \\
\bottomrule
\end{tabular}
\end{table}

\subsubsection{Discussão}
\todo{}

% =====================================================================
\section{Parte III --- Codec de áudio com perdas baseado na DCT}
% =====================================================================
\subsection{Visão geral e diagrama de blocos}
\todo{Figura com o diagrama de blocos.}

\subsection{Segmentação em blocos e janelamento}
\todo{Tamanho do bloco, sobreposição, janela; efeito nos artefactos de bloco.}

\subsection{Transformada DCT}
\todo{Tipo de DCT (DCT-II/IV, MDCT), implementação (p.\,ex.\ FFTW) e custo computacional.}

\subsection{Quantização dos coeficientes}
\subsubsection{Quantização uniforme}
\todo{}
\subsubsection{Quantização dependente da frequência}
\todo{Bandas, pesos, eliminação de coeficientes de alta frequência;
eventual modelo psicoacústico simplificado.}

\subsection{Codificação entrópica dos coeficientes}
\todo{Reutilização do Golomb da Parte II, codificação de zeros, etc.}

\subsection{Controlo de débito e opções do codificador}
\todo{Como o utilizador escolhe o compromisso débito--distorção.}

\subsection{Justificação experimental das decisões}
\todo{Curvas débito--distorção (SNR vs.\ kbps) para cada alternativa testada.}

\subsection{Comparação com codecs MPEG}
\todo{MP3 e AAC (e p.\,ex.\ Opus) via ffmpeg: taxa de compressão, SNR,
tempos e memória. Gráfico SNR vs.\ débito.}

\begin{figure}[h]
\centering
\fbox{\parbox[c][4cm][c]{0.7\linewidth}{\centering Gráfico SNR (dB) vs.\ débito (kbps)}}
% \includegraphics[width=0.7\linewidth]{figuras/rd_curve.pdf}
\caption{Desempenho débito--distorção do codec proposto e dos codecs MPEG.}
\label{fig:rd}
\end{figure}

\subsection{Discussão}
\todo{Limitações da SNR como medida de qualidade percetual, etc.}

% =====================================================================
\section{Conclusões e trabalho futuro}
% =====================================================================
\todo{Principais resultados e o que poderia ser melhorado.}

% =====================================================================
\section{Manual de utilização}
% =====================================================================
\subsection{Requisitos e instalação}
\todo{Dependências (apt), clonar o repositório, compilar.}

\subsection{Execução dos programas}
\todo{Sintaxe e exemplo para cada programa:
\prog{wav\_hist}, \prog{wav\_quant}, \prog{wav\_cmp}, \prog{wav\_effects},
\prog{wav\_lossless\_enc}/\prog{wav\_lossless\_dec},
\prog{wav\_lossy\_enc}/\prog{wav\_lossy\_dec}.}

\subsection{Reprodução dos resultados}
\todo{Scripts de benchmark e de geração de gráficos.}

% =====================================================================
\section*{Contribuição dos autores}
\addcontentsline{toc}{section}{Contribuição dos autores}
% =====================================================================
\todo{Percentagem e tarefas de cada membro (p.\,ex.\ 33/33/33).}

% =====================================================================
\begin{thebibliography}{9}
\bibitem{flac} Xiph.Org Foundation, \emph{FLAC --- Free Lossless Audio Codec}, \url{https://xiph.org/flac/}.
\bibitem{sndfile} E.~de Castro Lopo, \emph{libsndfile}, \url{https://libsndfile.github.io/libsndfile/}.
\bibitem{sayood} K.~Sayood, \emph{Introduction to Data Compression}, 5.\textsuperscript{a} ed., Morgan Kaufmann, 2017.
\end{thebibliography}

\end{document}
